% GENERATED FILE. Edit canonical lessons, then run npm run generate:books.
% canonical-source-hash: fnv1a64:afb655b09efa0a65
% canonical-lessons: KA-C475-aranya, KA-C475-malekadu, KA-C475-kadgiccu, KA-C475-kadi-fell, KA-C475-samraksisu, KA-R475-review-469-475

\chapter{Into the Woods}
\label{ch:ka-into-the-woods}
\hlchaptermodality{475}

\begin{chapteropening}
\textbf{By the end of this chapter:} \emph{I can say woodland, a rainforest, a forest fire, to fell and to conserve.}

Complete the last lesson of chapter 475: I can say woodland, a rainforest, a forest fire, to fell and to conserve.
\end{chapteropening}

\section[\texorpdfstring{\kn{ಅರಣ್ಯ} (araṇya)}{ಅರಣ್ಯ (araṇya)}]{\kn{ಅರಣ್ಯ} (araṇya) — woodland; a forest}

\label{lesson:KA-C475-aranya}



\begin{quote}
Before the new one: say the Kannada for to forsake, then the Kannada for solar power.
\end{quote}

\subsection*{You'll want to know: \kn{ಅರಣ್ಯ}}
\textbf{\kn{ಅರಣ್ಯ}} — \emph{araṇya} — \textquotedblleft{}woodland; a forest\textquotedblright{}.

\begin{quote}
\kn{ಅರಣ್ಯ} \kn{ಇಲಾಖೆ} \kn{ಹೊಸ} \kn{ಮರಗಳನ್ನು} \kn{ನೆಟ್ಟಿತು}.
\end{quote}

\begin{quote}
\emph{araṇya ilākhe hosa maragaḷannu neṭṭitu.}
\end{quote}

\begin{quote}
The forest department planted new trees.
\end{quote}

A formal word; \kn{ಕಾಡು} is the everyday forest.

\begin{cousinweb}
\textbf{\kn{ಅರಣ್ಯ}} (\emph{araṇya}) is borrowed from Sanskrit \emph{araṇya}, \textquotedblleft{}foreign land, wilderness, a forest\textquotedblright{} (CDIAL 600).
\end{cousinweb}

\begin{grammarlens}[title={What you've built}]
One more word for getting around.
\end{grammarlens}

\subsection*{Your turn}
\begin{itemize}
\raggedright
  \item \emph{Say it:} \emph{araṇya}
  \item \emph{Say it:} \emph{araṇya}, once more
  \item \emph{Say it:} \emph{araṇya ilākhe hosa maragaḷannu neṭṭitu.} — The forest department planted new trees.
  \item \emph{Recall:} say the Kannada for to forsake, then the Kannada for solar power, then say \emph{araṇya} again
\end{itemize}

\begin{tcolorbox}[breakable,colback=teal!4,colframe=teal!35!black,title={Before you move on}]
What does \kn{ಅರಣ್ಯ} mean? (\textquotedblleft{}Woodland; a forest\textquotedblright{}.) Say it once more.
\end{tcolorbox}
\section[\texorpdfstring{\kn{ಮಳೆಕಾಡು} (maḷekāḍu)}{ಮಳೆಕಾಡು (maḷekāḍu)}]{\kn{ಮಳೆಕಾಡು} (maḷekāḍu) — a rainforest}

\label{lesson:KA-C475-malekadu}



\begin{quote}
Before the new one: say the Kannada for solar power, then the Kannada for woodland.
\end{quote}

\subsection*{You'll want to know: \kn{ಮಳೆಕಾಡು}}
\textbf{\kn{ಮಳೆಕಾಡು}} — \emph{maḷekāḍu} — \textquotedblleft{}a rainforest\textquotedblright{}.

\begin{quote}
\kn{ಮಳೆಕಾಡಿನಲ್ಲಿ} \kn{ಸೂರ್ಯನ} \kn{ಬೆಳಕು} \kn{ನೆಲ} \kn{ಮುಟ್ಟುವುದಿಲ್ಲ}.
\end{quote}

\begin{quote}
\emph{maḷekāḍinalli sūryana beḷaku nela muṭṭuvudilla.}
\end{quote}

\begin{quote}
In a rainforest the sunlight does not reach the ground.
\end{quote}

\kn{ಮಳೆ}, rain, plus \kn{ಕಾಡು}, forest.

\begin{cousinweb}
\textbf{\kn{ಮಳೆಕಾಡು}} (\emph{maḷekāḍu}) is a native Kannada compound of \kn{ಮಳೆ}, \textquotedblleft{}rain\textquotedblright{} (DEDR 4753), and \kn{ಕಾಡು}, \textquotedblleft{}a forest\textquotedblright{} (DEDR 1438).
\end{cousinweb}

\begin{grammarlens}[title={What you've built}]
One more word for getting around.
\end{grammarlens}

\subsection*{Your turn}
\begin{itemize}
\raggedright
  \item \emph{Say it:} \emph{maḷekāḍu}
  \item \emph{Say it:} \emph{maḷekāḍu}, once more
  \item \emph{Say it:} \emph{maḷekāḍinalli sūryana beḷaku nela muṭṭuvudilla.} — In a rainforest the sunlight does not reach the ground.
  \item \emph{Recall:} say the Kannada for solar power, then the Kannada for woodland, then say \emph{maḷekāḍu} again
\end{itemize}

\begin{tcolorbox}[breakable,colback=teal!4,colframe=teal!35!black,title={Before you move on}]
What does \kn{ಮಳೆಕಾಡು} mean? (\textquotedblleft{}A rainforest\textquotedblright{}.) Say it once more.
\end{tcolorbox}
\section[\texorpdfstring{\kn{ಕಾಡ್ಗಿಚ್ಚು} (kāḍgiccu)}{ಕಾಡ್ಗಿಚ್ಚು (kāḍgiccu)}]{\kn{ಕಾಡ್ಗಿಚ್ಚು} (kāḍgiccu) — a forest fire}

\label{lesson:KA-C475-kadgiccu}



\begin{quote}
Before the new one: say the Kannada for woodland, then the Kannada for a rainforest.
\end{quote}

\subsection*{You'll want to know: \kn{ಕಾಡ್ಗಿಚ್ಚು}}
\textbf{\kn{ಕಾಡ್ಗಿಚ್ಚು}} — \emph{kāḍgiccu} — \textquotedblleft{}a forest fire\textquotedblright{}.

\begin{quote}
\kn{ಕಳೆದ} \kn{ಬೇಸಿಗೆಯಲ್ಲಿ} \kn{ಈ} \kn{ಬೆಟ್ಟದಲ್ಲಿ} \kn{ಕಾಡ್ಗಿಚ್ಚು} \kn{ಆಯಿತು}.
\end{quote}

\begin{quote}
\emph{kaḷeda bēsigeyalli ī beṭṭadalli kāḍgiccu āyitu.}
\end{quote}

\begin{quote}
Last summer there was a forest fire on this hill.
\end{quote}

\kn{ಕಾಡು}, forest, plus \kn{ಕಿಚ್ಚು}, fire: the k turns to g in the join.

\begin{cousinweb}
\textbf{\kn{ಕಾಡ್ಗಿಚ್ಚು}} (\emph{kāḍgiccu}) is a native Kannada compound of \kn{ಕಾಡು}, \textquotedblleft{}a forest\textquotedblright{} (DEDR 1438), and \emph{kiccu}, \textquotedblleft{}fire\textquotedblright{} (DEDR 1514).
\end{cousinweb}

\begin{grammarlens}[title={What you've built}]
One more word for getting around.
\end{grammarlens}

\subsection*{Your turn}
\begin{itemize}
\raggedright
  \item \emph{Say it:} \emph{kāḍgiccu}
  \item \emph{Say it:} \emph{kāḍgiccu}, once more
  \item \emph{Say it:} \emph{kaḷeda bēsigeyalli ī beṭṭadalli kāḍgiccu āyitu.} — Last summer there was a forest fire on this hill.
  \item \emph{Recall:} say the Kannada for woodland, then the Kannada for a rainforest, then say \emph{kāḍgiccu} again
\end{itemize}

\begin{tcolorbox}[breakable,colback=teal!4,colframe=teal!35!black,title={Before you move on}]
What does \kn{ಕಾಡ್ಗಿಚ್ಚು} mean? (\textquotedblleft{}A forest fire\textquotedblright{}.) Say it once more.
\end{tcolorbox}
\section[\texorpdfstring{\kn{ಕಡಿ} (kaḍi)}{ಕಡಿ (kaḍi)}]{\kn{ಕಡಿ} (kaḍi) — to fell; to chop}

\label{lesson:KA-C475-kadi-fell}



\begin{quote}
Before the new one: say the Kannada for a rainforest, then the Kannada for a forest fire.
\end{quote}

\subsection*{You'll want to know: \kn{ಕಡಿ}}
\textbf{\kn{ಕಡಿ}} — \emph{kaḍi} — \textquotedblleft{}to fell; to chop\textquotedblright{}.

\begin{quote}
\kn{ರಸ್ತೆ} \kn{ವಿಸ್ತರಿಸಲು} \kn{ಅವರು} \kn{ತುಂಬಾ} \kn{ಮರಗಳನ್ನು} \kn{ಕಡಿದರು}.
\end{quote}

\begin{quote}
\emph{raste vistarisalu avaru tumbā maragaḷannu kaḍidaru.}
\end{quote}

\begin{quote}
To widen the road they cut down a lot of trees.
\end{quote}

Past (I …): \kn{ಕಡಿದೆ}, \emph{kaḍide}.

\kn{ಮರ} \kn{ಕಡಿ}, to fell a tree; \kn{ಕತ್ತರಿಸು} is to cut with a blade.

\begin{cousinweb}
\textbf{\kn{ಕಡಿ}} (\emph{kaḍi}) is native Dravidian: DEDR 1125 lists Kannada \emph{kaḍi}, \textquotedblleft{}to cut, to chop, to fell\textquotedblright{}.
\end{cousinweb}

\begin{grammarlens}[title={What you've built}]
One more word for getting around.
\end{grammarlens}

\subsection*{Your turn}
\begin{itemize}
\raggedright
  \item \emph{Say it:} \emph{kaḍi}
  \item \emph{Say it:} \emph{kaḍi}, once more
  \item \emph{Say it:} \emph{raste vistarisalu avaru tumbā maragaḷannu kaḍidaru.} — To widen the road they cut down a lot of trees.
  \item \emph{Recall:} say the Kannada for a rainforest, then the Kannada for a forest fire, then say \emph{kaḍi} again
\end{itemize}

\begin{tcolorbox}[breakable,colback=teal!4,colframe=teal!35!black,title={Before you move on}]
What does \kn{ಕಡಿ} mean? (\textquotedblleft{}To fell; to chop\textquotedblright{}.) Say it once more.
\end{tcolorbox}
\section[\texorpdfstring{\kn{ಸಂರಕ್ಷಿಸು} (saṃrakṣisu)}{ಸಂರಕ್ಷಿಸು (saṃrakṣisu)}]{\kn{ಸಂರಕ್ಷಿಸು} (saṃrakṣisu) — to conserve; to protect}

\label{lesson:KA-C475-samraksisu}



\begin{quote}
Before the new one: say the Kannada for a forest fire, then the Kannada for to fell.
\end{quote}

\subsection*{You'll want to know: \kn{ಸಂರಕ್ಷಿಸು}}
\textbf{\kn{ಸಂರಕ್ಷಿಸು}} — \emph{saṃrakṣisu} — \textquotedblleft{}to conserve; to protect\textquotedblright{}.

\begin{quote}
\kn{ನಾವು} \kn{ಕಾಡುಗಳನ್ನು} \kn{ಸಂರಕ್ಷಿಸಬೇಕು}.
\end{quote}

\begin{quote}
\emph{nāvu kāḍugaḷannu saṃrakṣisabēku.}
\end{quote}

\begin{quote}
We must conserve our forests.
\end{quote}

Past (I …): \kn{ಸಂರಕ್ಷಿಸಿದೆ}, \emph{saṃrakṣiside}.

\kn{ಸಂರಕ್ಷಣೆ} is conservation; \kn{ಕಾಪಾಡು} is the everyday verb.

\begin{cousinweb}
\textbf{\kn{ಸಂರಕ್ಷಿಸು}} (\emph{saṃrakṣisu}) is built from Sanskrit \emph{saṃrakṣaṇa}, \textquotedblleft{}guarding, preservation, protection\textquotedblright{} (Monier-Williams), with the verb ending -\emph{isu}.
\end{cousinweb}

\begin{grammarlens}[title={What you've built}]
One more word for getting around.
\end{grammarlens}

\subsection*{Your turn}
\begin{itemize}
\raggedright
  \item \emph{Say it:} \emph{saṃrakṣisu}
  \item \emph{Say it:} \emph{saṃrakṣisu}, once more
  \item \emph{Say it:} \emph{nāvu kāḍugaḷannu saṃrakṣisabēku.} — We must conserve our forests.
  \item \emph{Recall:} say the Kannada for a forest fire, then the Kannada for to fell, then say \emph{saṃrakṣisu} again
\end{itemize}

\begin{tcolorbox}[breakable,colback=teal!4,colframe=teal!35!black,title={Before you move on}]
What does \kn{ಸಂರಕ್ಷಿಸು} mean? (\textquotedblleft{}To conserve; to protect\textquotedblright{}.) Say it once more.
\end{tcolorbox}
\section[(dialogue)]{Looking back: from Bags, Bakeries and Deliveries to Into the Woods}

\label{lesson:KA-R475-review-469-475}



\begin{quote}
Say the Kannada for a handbag and to swap.
\end{quote}

\subsection*{Your turn}
Say these aloud:

\begin{itemize}
\raggedright
  \item a handbag, to swap, a warranty, a bakery, a home delivery
  \item to widen, an entrance gate, a power cut, supplies, a telephone
  \item a din, to gather, to stare, to tug, to spill over
  \item to get damaged, to be off guard, to escape unharmed, shelter, a helpline
  \item the environment, pollution, to pollute, soot, carbon
  \item waste matter, recycling, to sort, to forsake, solar power
  \item woodland, a rainforest, a forest fire, to fell, to conserve
\end{itemize}

\begin{tcolorbox}[breakable,colback=teal!4,colframe=teal!35!black,title={Before you move on}]
A handbag? (\textbf{\kn{ಕೈಚೀಲ}}, \emph{kaicīla}.) To swap? (\textbf{\kn{ಅದಲುಬದಲು} \kn{ಮಾಡು}}, \emph{adalubadalu māḍu}.) A warranty? (\textbf{\kn{ವಾರಂಟಿ}}, \emph{vāraṇṭi}.) A bakery? (\textbf{\kn{ಬೇಕರಿ}}, \emph{bēkari}.) A home delivery? (\textbf{\kn{ಡೆಲಿವರಿ}}, \emph{ḍelivari}.) To widen? (\textbf{\kn{ವಿಸ್ತರಿಸು}}, \emph{vistarisu}.) An entrance gate? (\textbf{\kn{ಪ್ರವೇಶದ್ವಾರ}}, \emph{pravēśadvāra}.) A power cut? (\textbf{\kn{ಕಡಿತ}}, \emph{kaḍita}.) Supplies? (\textbf{\kn{ಸರಬರಾಜು}}, \emph{sarabarāju}.) A telephone? (\textbf{\kn{ದೂರವಾಣಿ}}, \emph{dūravāṇi}.) A din? (\textbf{\kn{ಗದ್ದಲ}}, \emph{gaddala}.) To gather? (\textbf{\kn{ಜಮಾಯಿಸು}}, \emph{jamāyisu}.) To stare? (\textbf{\kn{ದಿಟ್ಟಿಸು}}, \emph{diṭṭisu}.) To tug? (\textbf{\kn{ಜಗ್ಗು}}, \emph{jaggu}.) To spill over? (\textbf{\kn{ತುಳುಕು}}, \emph{tuḷuku}.) To get damaged? (\textbf{\kn{ಹಾಳಾಗು}}, \emph{hāḷāgu}.) To be off guard? (\textbf{\kn{ಏಮಾರು}}, \emph{ēmāru}.) To escape unharmed? (\textbf{\kn{ಪಾರಾಗು}}, \emph{pārāgu}.) Shelter? (\textbf{\kn{ಆಶ್ರಯ}}, \emph{āśraya}.) A helpline? (\textbf{\kn{ಸಹಾಯವಾಣಿ}}, \emph{sahāyavāṇi}.) The environment? (\textbf{\kn{ಪರಿಸರ}}, \emph{parisara}.) Pollution? (\textbf{\kn{ಮಾಲಿನ್ಯ}}, \emph{mālinya}.) To pollute? (\textbf{\kn{ಕಲುಷಿತಗೊಳಿಸು}}, \emph{kaluṣitagoḷisu}.) Soot? (\textbf{\kn{ಮಸಿ}}, \emph{masi}.) Carbon? (\textbf{\kn{ಇಂಗಾಲ}}, \emph{iṅgāla}.) Waste matter? (\textbf{\kn{ತ್ಯಾಜ್ಯ}}, \emph{tyājya}.) Recycling? (\textbf{\kn{ಮರುಬಳಕೆ}}, \emph{marubaḷake}.) To sort? (\textbf{\kn{ವಿಂಗಡಿಸು}}, \emph{viṅgaḍisu}.) To forsake? (\textbf{\kn{ತ್ಯಜಿಸು}}, \emph{tyajisu}.) Solar power? (\textbf{\kn{ಸೌರಶಕ್ತಿ}}, \emph{sauraśakti}.) Woodland? (\textbf{\kn{ಅರಣ್ಯ}}, \emph{araṇya}.) A rainforest? (\textbf{\kn{ಮಳೆಕಾಡು}}, \emph{maḷekāḍu}.) A forest fire? (\textbf{\kn{ಕಾಡ್ಗಿಚ್ಚು}}, \emph{kāḍgiccu}.) To fell? (\textbf{\kn{ಕಡಿ}}, \emph{kaḍi}.) To conserve? (\textbf{\kn{ಸಂರಕ್ಷಿಸು}}, \emph{saṃrakṣisu}.)
\end{tcolorbox}
